Remote sensing change detection has developed around four main deep-learning design families: Siamese CNN pipelines, attention-enhanced refinement networks, transformer-based temporal fusion models, and methods that make feature interaction more explicit through priors, flow, or matching. These lines of work address overlapping parts of the problem, but they differ sharply in temporal modeling, computational cost, and deployment practicality. This section positions \cosa{} within that landscape.

\subsection{CNN-Based Change Detection}
Foundational deep CD systems are largely Siamese CNN encoder--decoder models that process the two dates with shared weights and localize change through feature differencing or late fusion \cite{daudt2018fully,ronneberger2015unet,hongruixuan2019}. This design remains attractive because it is simple, efficient, and well aligned with pixel-wise CD supervision. Subsequent CNN variants improved decoder fusion and local representation quality through denser skip connections, deeper supervision, or lightweight multi-branch refinement, as seen in DSIFN \cite{zhang2020deeply}, SNUNet-CD \cite{fang2021snunetcd}, TinyCD \cite{codegoni2022tinycd}, LightCDNet \cite{10214556}, DASNet \cite{jie2020}, adaptive Siamese fusion \cite{yunzuo2025}, asymmetric semantic CD networks \cite{kunping2022}, and multitask Siamese designs \cite{xibing2025}.

These CNN-based methods establish strong practical baselines, but their temporal interaction is often dominated by fixed differencing or local convolutional fusion. As a result, long-range contextual reasoning, pseudo-change suppression, and precise boundary recovery remain difficult when changes are sparse, weak, or spatially ambiguous. The main limitation is therefore not feature extraction alone, but the lack of an explicit temporal control mechanism for decoder refinement. This limitation motivates \cosa{}, which strengthens temporal discrimination inside the decoder without abandoning the efficiency and simplicity of the Siamese CNN backbone.

\subsection{Attention-Based Methods}
Attention modules are a common extension to Siamese CNNs because they selectively reweight features that are likely to be useful for CD. STANet \cite{chen2020stanet} introduced spatial--temporal attention for remote sensing CD, while HANet \cite{10093022} and SemiSANet \cite{chengzhen2022} explored hierarchical and graph-attentive variants. Other works refine difference features with explicit attention branches, including AFDE-Net \cite{s2023}, ADDEDNet \cite{junheng2025}, and several multi-scale fusion architectures such as three-branch attention networks \cite{yan2024}, MDANet \cite{shanshan2024}, M2F2Net \cite{binhao2025}, wavelet-based aggregation \cite{jiangwei2024}, and lightweight multi-spectral attention models \cite{layth2025}.

These methods improve feature selection and often outperform plain CNN baselines, especially when multi-scale fusion is important. Their main limitation is that the attention mechanism is typically generic: it highlights informative channels or spatial regions, but it does not explicitly treat local bi-temporal correlation as a control variable for refinement. In addition, stacking multiple attention branches can erode the efficiency advantage of Siamese CNNs without fully resolving CD-specific issues such as pseudo-change suppression and stable precision/recall calibration. Accordingly, \cosa{} uses correlation-guided sampling to make temporal interaction explicit while keeping the refinement lightweight and locally focused.

\subsection{Transformer-Based Approaches}
Transformer-based CD models push temporal interaction further by using self-attention to capture long-range dependencies across the bi-temporal pair. Representative examples include BIT \cite{chen2022bit}, ChangeFormer \cite{bandara2022changeformer}, transformer-based Siamese CD \cite{w2022}, STransUNet \cite{jianghua2022}, multitask CNN--Transformer semantic CD \cite{wei2024}, and STeInFormer \cite{xiaowen2024}. More recent work has explored stronger pretraining and larger representation models, including bitemporal foundation-model integration \cite{chen2023time}, foundation-model-based remote sensing CD \cite{10438490}, and new large-scale evaluation settings such as JL1-CD \cite{liu2025jl1}. Recent reviews also show that transformer-style global fusion is now a major axis of CD research \cite{devansh2025,kamalakar2025}.

The trade-off is practical. Global attention improves contextual modeling, but it also increases memory and computational cost as image resolution grows, and many transformer solutions require backbone-level redesign rather than local refinement. Strong temporal fusion may also reduce the marginal value of later decoder enhancement, which is consistent with our near-neutral BIT result. The main limitation here is that stronger temporal modeling is often purchased through higher complexity and reduced plug-in compatibility. This limitation motivates \cosa{}, which targets adaptive temporal refinement at the decoder level without incurring the full cost of global all-pairs interaction.

\subsection{Correlation-Based and Feature Interaction Methods}
A smaller but important line of work tries to move beyond plain differencing by making cross-temporal interaction more explicit. Changer argues that feature interaction is central to CD and introduces dedicated interaction modeling \cite{fang2022changer}. Change Guiding Network injects a learned change prior to steer prediction \cite{10234560}, while SFSCDNet uses semantic flow to encode structured temporal correspondence \cite{k2024}. Related formulations also widen the supervision or interaction setting, for example by reframing object change detection with single-temporal supervision \cite{zheng2021change} or by integrating richer bi-temporal features through large pre-trained models \cite{chen2023time}.

These methods are closer to \cosa{} in spirit because they recognize that temporal correspondence should be modeled more deliberately. However, they generally treat interaction as another representation-learning block, change prior, or global fusion stage. They do not use local correlation as an explicit decoder-level control signal that decides where refinement should be strengthened or suppressed. They also rarely combine explicit temporal interaction with both multi-scale context aggregation and learnable residual calibration in one lightweight plug-in module. \cosa{} targets this gap by turning local correlation into an explicit decoder-side refinement signal and coupling it with multi-scale fusion and learnable residual control in a single plug-in design.

\subsection{Positioning \cosa{}}
\cosa{} is positioned between heavy backbone replacement and generic add-on attention. Relative to CNN baselines \cite{daudt2018fully,hongruixuan2019}, it injects stronger temporal discrimination without changing the encoder or prediction head. Relative to attention-based refinements \cite{chen2020stanet,10093022,yan2024}, it uses bi-temporal correlation explicitly rather than only reweighting features implicitly. Relative to transformer-style models \cite{chen2022bit,bandara2022changeformer,wei2024}, it keeps computation local and preserves plug-in compatibility. Relative to interaction- or prior-based designs \cite{fang2022changer,10234560,k2024}, it couples three coordinated mechanisms in the decoder: correlation-guided sampling, multi-scale fusion, and learnable residual gating.

That combination is the main distinction of \cosa{}. The module is designed for targeted decoder refinement rather than wholesale architectural replacement, and the ablation later in the paper shows that its gains do not come from any one component in isolation. Recent hybrid systems for specialized CD settings, such as SiamMask-ICDNet for island-building monitoring \cite{lebao2025}, further illustrate a broader trend: many gains come from scenario-specific redesign rather than from a general refinement block. \cosa{} instead targets the unresolved design gap of a lightweight mechanism that uses local temporal correlation as an explicit control signal while remaining deployable across heterogeneous Siamese backbones.
